% Einheit 4 von 4 – Nullstellen und Schnittpunkte berechnen (bank/quadratische-funktionen/e4.jsonl, zusammenbau v0.1)
%% TODO Zweigzeile Teil 1: was hier gelernt wird (Satz in Schülersprache aus dem Einheitstitel) – Einheit 4
\einheitenkopf[e4]{Einheit 4 von 4 · Nullstellen und Schnittpunkte berechnen}
\zweigzeile{ab Kl. 10 (am Gymnasium ab Kl. 9) · P10 oft}
\setcounter{aufgabe}{16}

%% TODO Titel: Ich-kann-Satz fehlt in der Bank (Platzhalter: Kettenname) – Nr. 17
%% TODO Anweisung: ein Satz über der ersten Teilaufgabe fehlt in der Bank – Nr. 17
\begin{aufgabe}{Nullstellen und Schnittpunkte}
%% TODO \rechenplatz{n}: Schrittzahl des Grundfalls fehlt in der Bank (nur bei fester Schreibform, 2.3 b)
\begin{teile}
\teil Parabel $f(x) = x^2 - 2x$ und Gerade $g(x) = x + 4$ – welche Gleichung setzt du an? Nur hinschreiben. \leerfeld = \leerfeld
\end{teile}
\begin{gleichungsraster}[2]
\gl{\text{$f(x) = (x - 1)^2 - 4$ – Nullstellen?}} &
\end{gleichungsraster}
\begin{teile}
\teil $f(x) = (x - 2)^2 + 3$ – wie viele Nullstellen? Begründe ohne Rechnung.
\end{teile}
\begin{gleichungsraster}[2]
\gl{\text{$f(x) = x^2 - 2x - 8$ – Nullstellen?}} &
\gl{\text{$f(x) = x^2 - 4x + 1$ – Nullstellen? Runde auf zwei Stellen.}} \\
\gl{\text{$f(x) = 3x^2 - 3x - 18$ – Nullstellen?}} &
\gl{\text{$f(x) = x^2 + x + 1$ – für welche $x$ ist $f(x) = 7$?}} \\
\gl{\text{$f(x) = x^2 + x - 4$ und $g(x) = 2x + 2$ – an welchen Stellen $x$ schneiden sich Parabel und Gerade?}} &
\gl{\text{$f(x) = (x - 2)^2 - 1$ und $g(x) = x - 3$ – an welchen Stellen $x$ schneiden sich Parabel und Gerade?}} \\
\gl{\text{$f(x) = x^2 - 4x + 1$ und $g(x) = 2x - 4$ – Schnittpunkte?}} &
\end{gleichungsraster}
\begin{teile}
\teil $f(x) = x^2 - 5$, $g(x) = 2x + 3$ – ist $P(4|11)$ ein Schnittpunkt? \janein
\teil $f(x) = (x - 2)^2 + 1$ – gib eine Gerade $g$ an, die mit der Parabel keinen Punkt gemeinsam hat. g(x) = \leerfeld
\end{teile}
\begin{gleichungsraster}[2]
\gl{\text{$f(x) = x^2 + 3x - 2$ und $h(x) = x^2 - x + 6$ – Schnittpunkt?}} &
\gl{\text{Die Parabel $f(x) = (x - 4)^2 - 6$ und die Gerade $g(x) = -2x + 5$ schneiden sich in zwei Punkten. Berechne ihre Koordinaten. (P10 2022 OS)}} \\
\gl{\text{Berechne die Schnittpunkte der Parabel $f(x) = x^2 - 10x + 22$ mit der $x$-Achse. Runde auf zwei Stellen. (P10 2025 OS)}} &
\gl{\text{Die Parabel $f(x) = -x^2 + 4x + 2$ ist nach unten geöffnet. Für welche $x$ ist $f(x) = -10$? (P10 2023 OS)}} \\
\end{gleichungsraster}
\end{aufgabe}

%% TODO Titel: Ich-kann-Satz fehlt in der Bank (Platzhalter: Kettenname) – Nr. 18
%% TODO Anweisung: ein Satz über der ersten Teilaufgabe fehlt in der Bank – Nr. 18
\begin{aufgabe}{Nullstellen und Schnittpunkte – Fehler finden, Begründen, Darstellungswechsel, Anwendung}
\begin{gleichungsraster}[2]
\gl{\text{Lina berechnet die Nullstellen von $f(x) = (x + 1)^2 - 25$: \rechnung{0 &= (x + 1)^2 - 25 \\ 25 &= (x + 1)^2 \\ x + 1 &= 5 \\ x &= 4} Finde den Fehler und korrigiere.}} &
\end{gleichungsraster}
\begin{teile}
\teil Warum liegt jede Nullstelle einer Parabel auf der $x$-Achse?
\teil $f(x) = (x - 2)^2 - 3$ und $g(x) = x - 3$ – lies die Schnittpunkte am Graphen ab und bestätige einen durch Einsetzen.

\begin{ksys}[xmin=-1,xmax=5,ymin=-4,ymax=2,ablesen]
\parabel{1}{2}{-3}{f}
\gerade{1}{-3}{g}
\end{ksys}
$S_1($ \leerfeld | \leerfeld ), $S_2($ \leerfeld | \leerfeld )
\end{teile}
\begin{gleichungsraster}[2]
\gl{\text{Ein Ball fliegt auf der Bahn $h(x) = -0{,}1(x - 6)^2 + 4{,}9$ ($x$: Entfernung in m, $h$: Höhe in m). Wie weit fliegt der Ball, bis er auf dem Boden aufkommt?}} & \\
\end{gleichungsraster}
\end{aufgabe}
